% ============================================================
% AX Hackathon — Phase 1 Blueprint
% Problem: Speech Disentanglement — Speaker-Specific Custom Word Detection
% Compile: xelatex blueprint.tex  (preferred) OR pdflatex blueprint.tex
% If metropolis theme is missing: tlmgr install beamertheme-metropolis pgfopts
% ============================================================
\documentclass[aspectratio=169,10pt]{beamer}

% ---------- Theme ----------
\usetheme{metropolis}
\metroset{
  progressbar=frametitle,
  numbering=fraction,
  block=fill,
  sectionpage=progressbar,
  subsectionpage=progressbar
}

% ---------- Packages ----------
\usepackage[utf8]{inputenc}
\usepackage{graphicx}
\usepackage{booktabs}
\usepackage{array}
\usepackage{tabularx}
\usepackage{xcolor}
\usepackage{amsmath,amssymb}
\usepackage{mathtools}
\usepackage{tikz}
\usepackage{fontawesome5}
\usetikzlibrary{positioning,arrows.meta,shapes.geometric,backgrounds,fit}

% ---------- Samsung-adjacent palette ----------
\definecolor{sriblue}{HTML}{1428A0}
\definecolor{sridark}{HTML}{0C2340}
\definecolor{sriaccent}{HTML}{1E88E5}
\definecolor{srired}{HTML}{B91C1C}
\definecolor{srigold}{HTML}{B45309}
\definecolor{srigreen}{HTML}{65A30D}
\definecolor{sribg}{HTML}{F8FAFC}
\definecolor{srigray}{HTML}{475569}

\setbeamercolor{normal text}{fg=sridark}
\setbeamercolor{frametitle}{bg=sridark,fg=white}
\setbeamercolor{progress bar}{fg=sriaccent,bg=sriaccent!30}
\setbeamercolor{title separator}{fg=sriaccent}
\setbeamercolor{alerted text}{fg=srired}

% ---------- Footer ----------
\setbeamertemplate{footline}{%
  \leavevmode\hbox{%
    \begin{beamercolorbox}[wd=.45\paperwidth,ht=2.4ex,dp=1ex,leftskip=1em]{author in head/foot}%
      \color{srigray}\small AX Hackathon Phase 1 Blueprint \textbar{} Team \textbf{MalayM09}%
    \end{beamercolorbox}%
    \begin{beamercolorbox}[wd=.4\paperwidth,ht=2.4ex,dp=1ex]{title in head/foot}%
      \color{srigray}\small Speech Disentanglement \textbar{} Speaker-Specific Custom Word Detection%
    \end{beamercolorbox}%
    \begin{beamercolorbox}[wd=.15\paperwidth,ht=2.4ex,dp=1ex,rightskip=1em]{date in head/foot}%
      \hfill\color{srigray}\small\insertframenumber\,/\,\inserttotalframenumber%
    \end{beamercolorbox}}%
}

% ---------- Helper commands ----------
\newcommand{\highlight}[1]{\textcolor{srired}{\textbf{#1}}}
\newcommand{\kpi}[1]{\textcolor{sriblue}{\textbf{#1}}}
\newcommand{\novel}[1]{\textcolor{srigold}{\textbf{#1}}}
\newcommand{\illus}[1]{illustrations/#1}

% ---------- Title info ----------
\title{Edge-AI Speech Disentanglement}
\subtitle{Speaker-Specific Custom Word Detection under -5\,dB SNR on <3\,M parameters}
\author{Team \textbf{MalayM09}  ·  Malay Mishra}
\institute{Samsung \texttt{ennovateX} AX Hackathon — Phase 1 Blueprint}
\date{April 2026}

% ============================================================
\begin{document}

% ============== SLIDE 1 — TEAM INTRODUCTION ==============
\begin{frame}[plain]
  \titlepage
  \vspace{-0.8em}
  \begin{center}
    \small\color{srigray}
    \textit{Apache-2.0 licensed submission \textbar{} All code, weights, and a derived open benchmark will be released publicly.}
  \end{center}
\end{frame}

\begin{frame}{Team Introduction}
  \vspace{-0.3em}
  \textbf{\color{sriblue}Team Name:} \texttt{MalayM09} \hfill
  \textbf{\color{sriblue}Problem:} \texttt{PS04} — Speech Disentanglement
  \vspace{1.2em}

  \begin{columns}[T, onlytextwidth]
    \column{0.07\textwidth}
    \column{0.62\textwidth}
      \begin{block}{Team Member  \small\color{srigray}(solo submission)}
        \renewcommand{\arraystretch}{1.18}
        \begin{tabular}{@{}ll@{}}
          Name     & Malay Mishra \\
          College  & BITS Pilani, Pilani campus \\
          Department & Computer Science \\
          Year     & 2026 \\
          Email ID & \texttt{F20220116@pilani.bits-pilani.ac.in} \\
        \end{tabular}
      \end{block}
    \column{0.31\textwidth}
      \small\color{srigray}\textit{Code, weights, and benchmarks released under \textbf{Apache-2.0} at \texttt{github.com/MalayM09/SRIB}.}
  \end{columns}
\end{frame}

% ============== SLIDE 2 — PROBLEM STATEMENT ==============
\begin{frame}{Problem Statement \& Why It Matters}
  \begin{columns}[T, onlytextwidth]
    \column{0.55\textwidth}
      \textbf{\color{sriblue}Selected problem.} Speech disentanglement for speaker-specific custom word detection at the extreme edge.
      \vspace{0.6em}

      \textbf{\color{sriblue}The challenge.}
      A voice trigger must wake a device \emph{only} when:
      \begin{itemize}\setlength\itemsep{2pt}
        \item the \highlight{right person} speaks
        \item the \highlight{right word} — rejecting phonetic look-alikes
        \item under crowd / babble / traffic noise (-5\,dB to 30\,dB SNR)
        \item at 0.5\,m to 5\,m distance
        \item with \highlight{$<$3\,M parameters} and \highlight{$<$0.2\,s xRT}
      \end{itemize}
      \vspace{0.4em}

      \textbf{\color{sriblue}Why it matters.} Enables private, offline-first, personalized wake-words on phones, earbuds, hearables, and smart-home hubs — without cloud offload or raw audio leaving the device.
    \column{0.45\textwidth}
      \begin{block}{Target KPIs (from brief)}
        \centering\small
        \begin{tabular}{@{}lr@{}}
          \toprule
          Metric & Target \\
          \midrule
          True Acceptance — Clean & $\geq 99\%$ \\
          True Acceptance — Noisy & $\geq 90\%$ \\
          False Acceptance        & $< 1$/hr \\
          SNR range               & -5 to 30\,dB \\
          Distance range          & 0.5 to 5\,m \\
          Speakers supported      & Male + Female \\
          Parameters              & $< 3$\,M \\
          Execution time (xRT)    & $< 0.2$\,s \\
          \bottomrule
        \end{tabular}
      \end{block}
      \vspace{0.3em}
      \small\color{srigray}
      \textit{Regime where standard log-Mel + ResNet pipelines break down and Transformer-based KWS misses the latency budget.}
  \end{columns}
\end{frame}

% ============== SLIDE 3 — PROPOSED SOLUTION (HIGH LEVEL) ==============
\begin{frame}{Proposed Solution — High Level}
  \textbf{\color{sriblue}One sentence.}
  A single $\sim$830\,K-parameter joint student (KWS backbone validated at \kpi{324\,K}) performs Keyword Spotting (KWS) and Speaker Verification (SV), distilled from two frozen SOTA teachers, with speaker-conditioned classification and a FAR-aware training loop.
  \vspace{0.6em}

  \begin{columns}[T, onlytextwidth]
    \column{0.33\textwidth}
      \begin{block}{\color{sriblue}1. Robust Frontend}
        Learnable PCEN on Mel — survives -5\,dB SNR where log-Mel collapses. Equivalent to LEAF in practice at a fraction of the compute.\footnotemark[1]
      \end{block}
    \column{0.33\textwidth}
      \begin{block}{\color{sriblue}2. Compact Joint Student}
        Modified BC-ResNet-8 with TRM branching, MQMHA pooling on SV head, and \novel{FiLM-conditioned KWS head} (speaker embedding modulates the classifier).
      \end{block}
    \column{0.33\textwidth}
      \begin{block}{\color{sriblue}3. Multi-Teacher KD + FAR Loop}
        WavLM-Base+ (phonetic) + ECAPA-TDNN (biometric) teachers. Training loop includes \novel{online $\tau_s$ calibration} that bakes the $<$1\,FA/hr operating point into the embedding geometry.
      \end{block}
  \end{columns}
  \vfill
  \footnotetext[1]{\color{srigray}\scriptsize LEAF analysis (arXiv:2404.06702): only PCEN sub-layer actually trains; filterbanks and low-pass stay near initialization.}
\end{frame}

% ============== SLIDE 4 — TECHNICAL DETAILS: ARCHITECTURE ==============
\begin{frame}{Technical Details — System Architecture}
  \vspace{-0.4em}
  \centering
  \includegraphics[width=\textwidth,height=0.72\textheight,keepaspectratio]{\illus{1_architecture.png}}
  \vspace{-0.3em}
  {\scriptsize\color{srigray}
   Shared trunk (blue) $\to$ Task Representation Modules (yellow) $\to$ KWS / SV branches (green) $\to$ FiLM modulation (red) conditions KWS on speaker embedding (stop-gradient on forward).}
\end{frame}

% ============== SLIDE 5 — TECHNICAL DETAILS: DISTILLATION ==============
\begin{frame}{Technical Details — Multi-Teacher Knowledge Distillation}
  \vspace{-0.4em}
  \centering
  \includegraphics[width=\textwidth,height=0.62\textheight,keepaspectratio]{\illus{2_distillation.png}}
  \vspace{-0.2em}
  {\scriptsize
  \[
    \mathcal{L}_{\text{total}} =
    \lambda_1\,\mathcal{L}_\text{KWS}^{\text{Focal}} +
    \lambda_2\,\mathcal{L}_\text{SV}^{\text{SubCtr-AAM}} +
    \lambda_3\,T^2\mathrm{KL}(\text{stu}_\text{kws}/T\,\|\,\text{probe}/T) +
    \lambda_4\,\mathcal{L}_\text{hint}^{\{4,8,11\}} +
    \lambda_5\,\mathcal{L}_\text{ECAPA}^{\text{multi-level}}
  \]}
\end{frame}

% ============== SLIDE 6 — TECHNICAL DETAILS: PARAM BUDGET & STAGES ==============
\begin{frame}{Technical Details — Parameter Budget \& Training Stages}
  \begin{columns}[T, onlytextwidth]
    \column{0.53\textwidth}
      \centering
      \includegraphics[width=\linewidth,height=0.75\textheight,keepaspectratio]{\illus{6_param_budget.png}}
    \column{0.47\textwidth}
      \small
      \begin{block}{4-Stage Training ($\sim$36 Kaggle GPU-hrs)}
        \begin{tabular}{@{}p{2.6cm}p{4.0cm}@{}}
          \toprule
          \textbf{Stage} & \textbf{Purpose} \\
          \midrule
          0. WavLM probe     & Train KWS probing head on frozen WavLM (one-time) \\
          1. Warm-up         & Stabilize PCEN \& trunk on SC only \\
          2. Joint MT + KD   & All 5 loss terms + FiLM + TRM active \\
          3. Trigger FT      & Custom-trigger specialization + \novel{online $\tau_s$ loop} \\
          4. QAT             & INT8 quant-aware training \\
          \bottomrule
        \end{tabular}
      \end{block}
      \textbf{\color{sriblue}Headroom:} 72\% under the 3\,M cap — room for QAT scale/zero-point tensors and late-stage capacity additions without refactor.
  \end{columns}
\end{frame}

% ============== SLIDE 7 — TIMELINE ==============
\begin{frame}{Technical Details — Execution Timeline}
  \vspace{-0.3em}
  \centering
  \includegraphics[width=\textwidth,height=0.65\textheight,keepaspectratio]{\illus{4_timeline.png}}
  \vspace{0.2em}
  \begin{columns}[T, onlytextwidth]
    \column{0.5\textwidth}
      \small\textbf{\color{sriblue}Compute plan:}
      Kaggle free tier (T4 / P100) — all training fits in 12-hour session limits with checkpoint-and-resume.
    \column{0.5\textwidth}
      \small\textbf{\color{sriblue}Deployment:}
      PyTorch QAT $\to$ ONNX opset 17 $\to$ ORT INT8 static. Benchmarked on Raspberry Pi 4B as conservative spec; also exportable to Galaxy NPU via QNN.
  \end{columns}
\end{frame}

% ============== SLIDE 8 — PHASE-1 EMPIRICAL VALIDATION ==============
\begin{frame}{Phase-1 Empirical Validation — KWS Backbone}
  \vspace{-0.3em}
  \begin{columns}[T, onlytextwidth]
    \column{0.62\textwidth}
      \includegraphics[width=\linewidth,height=0.70\textheight,keepaspectratio]{\illus{7_snr_sweep.png}}
      {\scriptsize\color{srigray}
       Noise-aug curves averaged over \{PCEN, log-Mel\} $\times$ \{seed 1337, seed 42\} = 4 runs; clean-training curve averaged over 2 front-ends. Panel (b) shaded bands = halfspread over 2 seeds. Test-time MUSAN mixed at fixed SNR on the val split (identical noise draw across arms for a given seed).}
    \column{0.38\textwidth}
      \small
      \begin{block}{\color{sriblue}Headline}
        \kpi{97.6\%} clean val on SC V2 @ \kpi{324\,K} params (10.8\% of the 3\,M cap). Train/val gap $<$2\,pp — no overfit.
      \end{block}
      \begin{block}{\color{srired}Dominant robustness lever}
        \textbf{Data diversity, not the front-end.}\\
        At worst case (RIR + $-$5\,dB noise): noise-aug training holds \kpi{77\%}; clean training collapses to \alert{55\%}. \highlight{$+$22\,pp gap} from augmentation alone.
      \end{block}
      \begin{block}{\color{sriblue}PCEN vs log-Mel (under noise-aug)}
        Direction \emph{stable} across 2 seeds at every SNR — PCEN ahead by mean +0.41\,pp. PCEN also shows \kpi{\~4$\times$ lower cross-seed variance} at low SNR — a reproducibility win, not just an accuracy win.
      \end{block}
  \end{columns}
\end{frame}

% ============== SLIDE 8b — SV BRANCH POC ==============
\begin{frame}{Phase-1 Empirical Validation — SV Branch POC}
  \vspace{-0.3em}
  \begin{columns}[T, onlytextwidth]
    \column{0.60\textwidth}
      \includegraphics[width=\linewidth,height=0.70\textheight,keepaspectratio]{\illus{8_sv_progression.png}}
      {\scriptsize\color{srigray}
       Trained on VoxCeleb1 dev (1211 speakers, 24K utterances), evaluated on the canonical VoxCeleb1-O trial list (37,720 same/different pairs). Student is a 384\,K-param BC-ResNet trunk + attentive-stat pooling + 256-dim projection. Teacher: pretrained \texttt{resemblyzer} (GE2E-LSTM).}
    \column{0.40\textwidth}
      \small
      \begin{block}{\color{sriblue}Current state}
        \kpi{10.64\,\%} EER on VoxCeleb1-O after 40 epochs joint distillation + Sub-center AAM-Softmax. Student is 384\,K params (within budget for joint student).
      \end{block}
      \begin{block}{\color{srired}Why v1 collapsed at 28.67\,\%}
        Pure cosine distillation with no classification signal $\Rightarrow$ embeddings cluster at a single point — \emph{no} angular discrimination. Sub-center AAM with margin $m{=}0.2$ forces speakers into separate angular cones; \highlight{$+$18\,pp gain} from the loss change alone.
      \end{block}
      \begin{block}{\color{srigold}Path to $\leq$5\,\% (v3)}
        Three remaining levers: (i) swap teacher to ECAPA-TDNN (1\,\% EER vs 4\,\%); (ii) tune $\lambda_{\text{distill}}$ down so the two losses cooperate; (iii) train on full VoxCeleb1 ($\sim$150K utts) for 80 epochs. Projected v3 EER: \kpi{4--5\,\%} — competitive with x-vector at $\sim$10$\times$ smaller.
      \end{block}
  \end{columns}
\end{frame}

% ============== SLIDE 8c — FiLM SPEAKER-CONDITIONED REJECTION ==============
\begin{frame}{Phase-1 Empirical Validation — FiLM Rejection Trade-off}
  \vspace{-0.3em}
  \begin{columns}[T, onlytextwidth]
    \column{0.62\textwidth}
      \includegraphics[width=\linewidth,height=0.68\textheight,keepaspectratio]{\illus{9_film_tradeoff.png}}
      {\scriptsize\color{srigray}
       Joint KWS + SV + FiLM model (500\,K params, 400 train speakers, 8K utts, 30 epochs). Eval on 28 held-out speakers $\times$ 10 test items each, with 5-clip enrollment averaging. Genuine = enrolled speaker says keyword; imposter = different speaker enrolled, keyword spoken by another speaker. Collapse ratio = mean P(true\_kw)$_\text{imposter}$ / mean P(true\_kw)$_\text{genuine}$.}
    \column{0.38\textwidth}
      \small
      \begin{block}{\color{sriblue}Mechanism validated}
        \kpi{95\,\% TAR + 86\,\% imposter-rejection} at \kpi{14\,\% FAR} when train and eval embedding distributions are matched (v2 own-audio diagnostic). Probability collapse ratio \kpi{0.20} — imposters get $\sim$5$\times$ less keyword confidence than genuine.
      \end{block}
      \begin{block}{\color{srired}Trade-off characterized}
        Closing the train/eval gap (v3 paired enrollment) lifts TAR to 92\,\% but loosens rejection (FAR jumps to \alert{48\,\%}). The system has a tunable TAR/FAR knob via imposter-loss weighting and hard-negative mining.
      \end{block}
      \begin{block}{\color{srigold}Path to deployment (v4)}
        Hard-negative mining + $\lambda_{\text{imp}}{=}2$ + max-imposter-ratio 0.7. Target: \kpi{$\geq$85\,\% TAR + $\leq$25\,\% FAR + $\geq$70\,\% REJ} at the production operating point.
      \end{block}
  \end{columns}
\end{frame}

% ============== SLIDE 8d — CUSTOM-WORD FEW-SHOT VALIDATION ==============
\begin{frame}{Phase-1 Empirical Validation — Custom-Word Generalization}
  \vspace{-0.3em}
  \begin{columns}[T, onlytextwidth]
    \column{0.62\textwidth}
      \includegraphics[width=\linewidth,height=0.68\textheight,keepaspectratio]{\illus{10_custom_word.png}}
      {\scriptsize\color{srigray}
       Embedding network (384\,K params) trained on 28 SC V2 keywords with Sub-center AAM-Softmax. Evaluated on 7 held-out polysyllabic words (\texttt{backward, forward, marvin, sheila, learn, follow, visual}) the model has \emph{never seen}: 5-shot enrollment averaging into a prototype, then cosine similarity scoring against genuine + cross-word imposter pairs (700 trials).}
    \column{0.38\textwidth}
      \small
      \begin{block}{\color{sriblue}Headline}
        \kpi{EER = 12.43\,\%} on novel words, \kpi{AUC = 0.934}, score separation \kpi{4.2$\times$} (genuine 0.63 vs imposter 0.15) — at \kpi{384\,K params, no WavLM teacher}.
      \end{block}
      \begin{block}{\color{srired}Phonetic generalization confirmed}
        Per-word EER spans 6--18\,\%: words sharing phonemes with the train set (\texttt{forward}, \texttt{follow}) enroll at \kpi{6--7\,\%}; novel phoneme combos (\texttt{backward}, \texttt{learn}) at \alert{16--18\,\%}. The model is \emph{phonetically organized}, not memorising — a precondition for arbitrary user-chosen triggers.
      \end{block}
      \begin{block}{\color{srigold}Path to GREEN ($\leq$10\,\% EER)}
        Add WavLM hint-layer distillation (Slide 5 already specifies layers 4/8/11) — closes the 6--18\,\% spread on hard words. Projected: \kpi{$\leq$8\,\% EER, TAR$\geq$70\,\% at FAR=1\,\%}. Same 384\,K student, no extra inference cost.
      \end{block}
  \end{columns}
\end{frame}

% ============== SLIDE 8e — DEPLOYMENT VALIDATION ==============
\begin{frame}{Phase-1 Empirical Validation — Deployment Latency \& Size}
  \vspace{-0.3em}
  \begin{columns}[T, onlytextwidth]
    \column{0.62\textwidth}
      \includegraphics[width=\linewidth,height=0.68\textheight,keepaspectratio]{\illus{11_deploy_metrics.png}}
      {\scriptsize\color{srigray}
       Joint KWS+SV+FiLM model (500\,K params) exported to ONNX. Frontend (mel + PCEN) runs in PyTorch CPU; trunk + heads run in ONNX Runtime, single-thread (matches Pi 4B per-core profile). Mac M1 latencies are direct measurements (200 iters); Pi 4B numbers projected via the documented 3.5$\times$ M1$\to$Cortex-A72 single-core ratio. INT8 via ONNX static post-training quantization with 64 SC V2 calibration clips.}
    \column{0.38\textwidth}
      \small
      \begin{block}{\color{sriblue}xRT  ·  comfortable PASS}
        Pi 4B INT8: \kpi{12\,ms} end-to-end vs \kpi{200\,ms} budget — \kpi{17$\times$ headroom}. Even FP32 at 17\,ms passes the budget by 12$\times$. The xRT KPI is not the binding constraint.
      \end{block}
      \begin{block}{\color{sridark}Param + size budget}
        \kpi{499.6\,K params} ($\sim$\kpi{17\%} of the 3\,M cap), INT8 ONNX ships at \kpi{0.79\,MB} (2.6$\times$ shrink from 2.05\,MB FP32). Tiny enough for OTA updates and on-device caching.
      \end{block}
      \begin{block}{\color{srigold}Numeric fidelity}
        FP32 ONNX: max-abs-diff vs PyTorch \kpi{$<$3$\times$10$^{-6}$} (essentially exact). INT8 post-hoc: KWS-logit drift $\sim$2.0, SV-emb drift $\sim$0.16 — expected for static PTQ; production deployment uses \novel{INT8 QAT} (Stage 4 in the training plan), which typically reduces drift 5--10$\times$.
      \end{block}
  \end{columns}
\end{frame}

% ============== SLIDE 8f — BABBLE / COCKTAIL-PARTY ROBUSTNESS ==============
\begin{frame}{Phase-1 Empirical Validation — Cocktail-Party Robustness}
  \vspace{-0.3em}
  \begin{columns}[T, onlytextwidth]
    \column{0.62\textwidth}
      \includegraphics[width=\linewidth,height=0.68\textheight,keepaspectratio]{\illus{12_babble_robustness.png}}
      {\scriptsize\color{srigray}
       Babble noise = N random SC V2 utterances mixed at fixed SNR (cocktail-party simulation, addresses the \textit{``background conversations''} phrase in the PS04 brief). Models: P0 = clean training, P1a = RIR + MUSAN training, P1c = RIR + MUSAN + \novel{babble} training. Same val set, same noise seed across runs.}
    \column{0.38\textwidth}
      \small
      \begin{block}{\color{sriblue}Cocktail-party gap, identified}
        Without babble training, the model collapses on competing-speaker noise: P1a hits only \alert{29\,\%} at $-$5\,dB babble vs \kpi{84\,\%} at $-$5\,dB MUSAN. \kpi{$-$56\,pp gap} between ambient noise and speech-on-speech.
      \end{block}
      \begin{block}{\color{srigold}Babble augmentation, validated}
        Adding babble to training (40\,\% of batches, 1--2 competing voices) lifts babble robustness by \kpi{+19\,pp at $-$5\,dB}, \kpi{+21\,pp at 0\,dB}, with negligible clean cost ($-$0.5\,pp). N=3 generalization holds: 90\,\% at 10\,dB even with 3 simultaneous voices.
      \end{block}
      \begin{block}{\color{srired}Trade-off characterized}
        MUSAN robustness erodes \kpi{$\leq$5\,pp} at $-$5\,dB (84\,\% $\to$ 80\,\%) — within bounds. \novel{Net trade is strongly positive}: +19\,pp babble for $-$5\,pp MUSAN.
      \end{block}
  \end{columns}
\end{frame}

% ============== SLIDE 9 — NOVELTY HERO: FAR-IN-THE-LOOP ==============
\begin{frame}{Novelty \& Innovation — FAR-in-the-Loop Training}
  \vspace{-0.3em}
  \begin{columns}[T, onlytextwidth]
    \column{0.62\textwidth}
      \includegraphics[width=\linewidth,height=0.68\textheight,keepaspectratio]{\illus{3_far_in_loop.png}}
    \column{0.38\textwidth}
      \small
      \textbf{\color{srired}Core invention.}
      The deployment false-accept threshold $\tau_s$ is re-estimated every 500 training steps on a held-out imposter pool. The resulting FRR-at-target-FAR is fed back as a \novel{differentiable soft penalty}.
      \vspace{0.5em}

      \textbf{\color{sriblue}Why it matters.}
      Conventional pipelines treat $\tau_s$ as a post-hoc deployment knob. If the training loss never sees the $<$1\,FA/hr constraint, the embedding geometry cannot cluster tightly enough for \emph{any} threshold to satisfy it.\footnotemark[2]
      \vspace{0.5em}

      \textbf{\color{sriblue}Result.}
      The operating point is baked into the embedding geometry — not tuned around it.
  \end{columns}
  \footnotetext[2]{\color{srigray}\scriptsize PK-MTL (arXiv:2206.13708) reports vanilla BC-ResNet hits 64.32\% FRR in target-only KWS; task-adaptive training brings this to 6.56\%.}
\end{frame}

% ============== SLIDE 10 — NOVELTY: ADDITIONAL INNOVATIONS ==============
\begin{frame}{Novelty \& Innovation — Supporting Contributions}
  \small
  \begin{columns}[T, onlytextwidth]
    \column{0.5\textwidth}
      \begin{block}{\color{srired}Structural FA rejection via FiLM}
        KWS classifier weights are \novel{speaker-conditioned}. Imposter audio flows through the wrong $(\gamma,\beta)$ modulation and the keyword logit collapses — structural rejection, not threshold-based.\footnotemark[3]
      \end{block}
      \begin{block}{\color{srired}Layer-ensemble WavLM distillation}
        Learnable softmax weights over WavLM layers $\{4, 8, 11\}$ — outperforms single-layer distillation.\footnotemark[4]
      \end{block}
      \begin{block}{\color{srired}Multi-level ECAPA distillation}
        Cosine on final embedding + MSE on block-2 and block-3 features (projected via $1{\times}1$ conv). $\sim$5\% relative EER gain.\footnotemark[5]
      \end{block}
    \column{0.5\textwidth}
      \begin{block}{Considered \& rejected alternatives}
        \renewcommand{\arraystretch}{1.1}
        \begin{tabular}{@{}p{2.3cm}p{5.0cm}@{}}
          \toprule
          \textbf{Option} & \textbf{Rejected because} \\
          \midrule
          Audio Mamba & Needs 5$\times$ BC-ResNet params for equal clean acc.; no noise-robust results; deployment immature \\
          LEAF / SincNet & LEAF analysis shows only PCEN actually trains — already covered \\
          GRL adversarial & Removes the speaker signal we \emph{need} for personalized KWS \\
          Deep cross-attn & 4$\times$ compute — breaks xRT budget \\
          \bottomrule
        \end{tabular}
      \end{block}
  \end{columns}
  \footnotetext[3]{\color{srigray}\scriptsize FiLM for personalized KWS: arXiv:2311.03419 — 2-6\% relative EER reduction at $<$1\% param overhead.}
  \footnotetext[4]{\color{srigray}\scriptsize SKILL, Interspeech 2024 (arXiv:2402.16830).}
  \footnotetext[5]{\color{srigray}\scriptsize Multi-level X-vector KD, arXiv:2303.01125.}
\end{frame}

% ============== SLIDE 10 — ENROLLMENT + DEMO PLAN ==============
\begin{frame}{Deployment \& Demo Plan (for Grand Finale)}
  \vspace{-0.2em}
  \begin{columns}[T, onlytextwidth]
    \column{0.56\textwidth}
      \includegraphics[width=\linewidth,height=0.78\textheight,keepaspectratio]{\illus{5_enrollment_runtime.png}}
    \column{0.44\textwidth}
      \small
      \begin{block}{Live demo script}
        \begin{enumerate}\setlength\itemsep{3pt}
          \item Enroll a live judge: 5 clean utterances ($\sim$30\,s).
          \item \textbf{Clean trigger} from judge $\to$ \kpi{fires}.
          \item \textbf{Noisy trigger} — traffic at -5\,dB from phone speaker $\to$ \kpi{fires}.
          \item \textbf{Imposter} (teammate) says the same word $\to$ \alert{rejected}.
          \item \textbf{Confusable} phrase (`turn on' vs `turn off') $\to$ \alert{rejected}.
        \end{enumerate}
      \end{block}
      \begin{block}{Hardware}
        Raspberry Pi 4B (conservative) or Galaxy-class NPU via QNN. No network connection during demo.
      \end{block}
  \end{columns}
\end{frame}

% ============== SLIDE 11 — OPEN DATASETS ==============
\begin{frame}{Open Datasets — Used \& Published}
  \small
  \begin{block}{To be used (all permissively licensed)}
    \centering
    \begin{tabular}{@{}p{3.5cm}p{7.5cm}p{2.0cm}@{}}
      \toprule
      \textbf{Dataset} & \textbf{Role} & \textbf{Licence} \\
      \midrule
      Google Speech Commands v0.02 & KWS positives + general negatives & CC-BY 4.0 \\
      VoxCeleb1                    & Speaker verification + imposter pool & CC-BY 4.0 \\
      LibriPhrase                  & Joint phrase + speaker fine-tuning & MIT \\
      MUSAN                        & Crowd / babble / traffic noise injection (-5 to 30\,dB) & CC-BY 4.0 \\
      OpenSLR RIRs                 & Room impulse responses — 0.5--5\,m distance simulation & Apache-2.0 \\
      CMUdict + TIMIT phone-conf.  & G2P + phonetic confusable mining & BSD / public \\
      \bottomrule
    \end{tabular}
  \end{block}
  \begin{block}{\color{srigold}To be published (open-source deliverable)}
    \textbf{Confusable-Trigger Benchmark.} For any chosen trigger word: the top-50 mined phonetic confusables synthesized in 100+ XTTS-v2 voices, each passed through the RIR + MUSAN augmentation pipeline. \novel{First public benchmark explicitly targeting the $<$1\,FA/hr regime} in personalized KWS. Released under \textbf{Apache-2.0}.
  \end{block}
\end{frame}

% ============== SLIDE 12 — OPEN MODELS ==============
\begin{frame}{Open Models — Used \& Developed}
  \small
  \begin{block}{Used as frozen teachers / utilities}
    \centering
    \begin{tabular}{@{}p{3.2cm}p{6.0cm}p{3.5cm}@{}}
      \toprule
      \textbf{Model} & \textbf{Role} & \textbf{Source / Licence} \\
      \midrule
      WavLM-Base+  & Phonetic, noise-robust teacher    & \texttt{microsoft/wavlm-base-plus} \ MIT \\
      ECAPA-TDNN   & Speaker biometric teacher          & \texttt{speechbrain/spkrec-ecapa} \ Apache-2.0 \\
      XTTS-v2      & Multi-voice TTS for trigger / confusable augmentation & \texttt{coqui/XTTS-v2} \ CPML \\
      g2p-en       & Trigger G2P $\to$ ARPAbet phones   & Apache-2.0 \\
      \bottomrule
    \end{tabular}
  \end{block}
  \begin{block}{\color{srigold}Developed \& trained from scratch (our contribution)}
    \centering
    \begin{tabular}{@{}p{6.0cm}p{4.5cm}p{2.0cm}@{}}
      \toprule
      \textbf{Model} & \textbf{Description} & \textbf{Licence} \\
      \midrule
      Modified BC-ResNet-8 + TRM + FiLM + MQMHA  & Joint KWS + SV student, $\sim$830\,K params (trunk validated at 324\,K) & Apache-2.0 \\
      WavLM $\to$ KWS probing head                & 2-layer MLP, $\sim$50\,K params              & Apache-2.0 \\
      INT8 ONNX deployable artifact               & Streaming, stateful PCEN                     & Apache-2.0 \\
      \bottomrule
    \end{tabular}
  \end{block}
\end{frame}

% ============== SLIDE 13 — AI / GENAI / TOOLS & BEST PRACTICES ==============
\begin{frame}{AI / GenAI / Agentic Tools \& Best Practices Uncovered}
  \small
  \begin{columns}[T, onlytextwidth]
    \column{0.45\textwidth}
      \begin{block}{\color{sriblue}Tools}
        \begin{itemize}\setlength\itemsep{2pt}
          \item \textbf{PyTorch 2.x} + \texttt{torch.compile} — $\sim$15\% train speedup.
          \item \textbf{HuggingFace / SpeechBrain} for frozen teacher inference.
          \item \textbf{Coqui XTTS-v2} — generative-AI speaker diversity for trigger synthesis.
          \item \textbf{LLM-driven adversarial confusable generator} (Llama-3-8B, in-context prompt) extends CMUdict-based mining with open-vocabulary near-misses.
          \item \textbf{Optuna TPE} — automated Stage-2 multi-task hyperparameter search.
          \item \textbf{ONNX Runtime + Quantization Toolkit} — edge export + benchmarking.
          \item \textbf{Google \texttt{kws\_streaming}} — reference for stateful streaming inference.
        \end{itemize}
      \end{block}
    \column{0.55\textwidth}
      \begin{block}{\color{srigold}Best practices documented in release}
        \begin{enumerate}\setlength\itemsep{2pt}
          \item \textbf{Symmetric augmentation in distillation.} Teacher and student must see the \emph{same} augmented audio — asymmetric pairing collapses the student to the teacher's clean manifold.
          \item \textbf{Multi-level KD beats single-layer KD} for both phonetic (WavLM) and biometric (ECAPA) teachers.
          \item \textbf{Bake the operating point into training.} Online $\tau_s$ calibration — FAR-in-the-Loop.
          \item \textbf{Structural conditioning $>$ threshold gating.} FiLM-modulated KWS outperforms the AND-of-two-classifiers pattern on FA.
          \item \textbf{Curriculum on SNR, not just on data.} Ramp 30\,dB $\to$ -5\,dB across first 15 epochs.
          \item \textbf{TRM $>$ PCGrad} for MTL gradient decoupling in this regime.
        \end{enumerate}
      \end{block}
  \end{columns}
  \vspace{0.2em}
  \centering
  \small\color{srigray}\textit{Full licence statement: all code, model weights, and benchmarks released under Apache-2.0, per hackathon rules.}
\end{frame}

% ============== CLOSING ==============
\begin{frame}[standout]
  \centering
  \textbf{Closing the FA $<$ 1/hr gap on $<$3\,M params}\\[0.5em]
  \normalsize FAR-in-the-Loop training  $\cdot$  FiLM-conditioned KWS  $\cdot$  Multi-teacher distillation\\[1em]
  \small\color{srigray}Team \texttt{MalayM09}  $\cdot$  Malay Mishra  $\cdot$  BITS Pilani  $\cdot$  Problem \texttt{PS04}  $\cdot$  Apache-2.0
\end{frame}

\end{document}
